\chapter{Robotersteuerung}
\label{ch:robotersteuerung}
% Verantwortlich: Person 2

Die Robotersteuerung setzt die Bewegungsaufträge des Orchestrators am Roboter um: Sie greift
das Bauteil an der erkannten Pose, legt es an der Prüfposition ab, nimmt es dort nach der
Prüfung wieder auf und legt es auf der Ablage für Gut- bzw.\ Schlechtteile ab. Nach außen
bietet sie dafür die Schnittstelle \texttt{RobotController} aus \cref{sec:schnittstellen}
an. Kenntnis über Kameras oder Prüfergebnisse benötigt sie nicht.

% ------------------------------------------------------------------------------
\section{Roboter und Programmierschnittstelle}
\label{sec:m2-api}

Eingesetzt wird ein Neura Robotics LARA, ein kollaborativer Roboter (\ac{cobot}) mit sechs
Achsen (\cref{sec:robotik}). Die Steuerung läuft mit der Softwareversion~5.0.8. Für die
Ansteuerung von einem externen Rechner stellt der Hersteller die Python-Bibliothek NeuraPy
bereit~\cite{neura2025neurapy}. Sie sendet Befehle als JSON-Nachrichten über eine
TCP-Verbindung an die Steuerung und erhält deren Antworten auf demselben Weg.

Kartesische Posen werden in NeuraPy als $[x, y, z, \text{roll}, \text{pitch}, \text{yaw}]$
in Metern und Radiant angegeben. Innerhalb des Projekts werden dagegen durchgängig
Millimeter und Grad verwendet (\cref{sec:schnittstellen}); die Umrechnung erfolgt
ausschließlich in der Robotersteuerung.

Bevor der Roboter Bewegungen ausführt, wird er vorbereitet: Die Antriebe werden
eingeschaltet, die Steuerung wird bei Bedarf vom Handbetrieb in den Automatikbetrieb
umgeschaltet, das Werkzeug des Greifers wird ausgewählt, und Geschwindigkeitsbegrenzungen
werden gesetzt. Abschließend wird auf der Steuerung ein Bewegungsprogramm gestartet, ohne das
keine Bewegungsbefehle angenommen werden. Diese Vorbereitung läuft beim Verbinden und nach
jedem Stopp erneut ab.

Für die Bewegungen werden zwei Bewegungsarten genutzt (\cref{sec:robotik}):
\begin{itemize}
  \item \textbf{\ac{ptp}-Bewegung} für die großen Wege zwischen den Stationen. Aus der
        kartesischen Zielpose werden mit der inversen Kinematik der Steuerung die
        Gelenkwinkel berechnet, die dann in einer Gelenkbewegung angefahren werden.
  \item \textbf{\ac{lin}-Bewegung} für das Absenken auf das Bauteil und das Abheben. Der
        \ac{tcp} bewegt sich dabei auf einer Geraden, sodass der Greifer senkrecht auf das
        Bauteil zufährt.
\end{itemize}

% ------------------------------------------------------------------------------
\section{Greifpose}
\label{sec:m2-greifpose}

Aus der \texttt{ObjectPose} der Objekterkennung wird eine Zielpose für den Greifer
abgeleitet. Die Position übernimmt $x$ und $y$ direkt; die Höhe ergibt sich aus der
Tischhöhe plus einer festen Greifhöhe, die von Bauteil und Greifer abhängt. Der Greifer zeigt
immer senkrecht nach unten. Dazu wird die Orientierung der Ruheposition übernommen und nur um
die Hochachse des Roboters gedreht, und zwar um den erkannten Bauteilwinkel $\theta$ plus
einen festen Greifer-Offset. Der Offset von \ang{90} sorgt dafür, dass sich die Greiferbacken
quer zur Längsachse schließen, also an den beiden langen Seiten des Bauteils anliegen. Der
resultierende Winkel wird auf den Bereich \ang{-180} bis \ang{180} normiert.

\todo[inline]{Greifer-Offset und Drehrichtung der Hochachse bei nach unten zeigendem Werkzeug
(roll = \ang{180}) am realen Roboter prüfen.}

% ------------------------------------------------------------------------------
\section{Greif- und Ablegesequenz}
\label{sec:m2-sequenz}

Greifen und Ablegen laufen nach demselben Muster ab: Der Roboter fährt zunächst eine
Vorposition oberhalb des Ziels an, senkt sich dann langsam und geradlinig ab, schließt bzw.\
öffnet den Greifer und hebt wieder auf die Vorposition an (\cref{fig:m2-pick}). Die
Vorposition liegt einen festen Sicherheitsabstand über dem Ziel, sodass der Roboter bei der
schnelleren Gelenkbewegung nicht mit dem Bauteil oder dem Tisch kollidiert. Nach dem Schließen
und Öffnen des Greifers wird eine kurze Zeit gewartet, bis der Greifer seine Endlage erreicht
hat.

\begin{figure}[htbp]
  \centering
  \begin{minipage}[c]{0.44\textwidth}
    \centering
    \begin{tikzpicture}[font=\small,
        stepnr/.style={circle, draw, fill=white, inner sep=0.8pt, font=\scriptsize\bfseries}]
      \fill[black!10] (0,-0.25) rectangle (6.4,0);
      \draw[thick] (0,0) -- (6.4,0);
      \node[font=\scriptsize, anchor=north west] at (0,-0.25) {Tisch};
      \draw[fill=violet!25] (4.3,0) rectangle (5.3,0.5);
      \draw[fill=black!45] (4.13,0.12) rectangle (4.3,0.8);
      \draw[fill=black!45] (5.3,0.12) rectangle (5.47,0.8);
      \draw[fill=black!25] (4.13,0.8) rectangle (5.47,1.0);
      \coordinate (S) at (0.6,3.3);
      \coordinate (V) at (4.8,2.7);
      \fill (S) circle (1.6pt) node[below, inner sep=3pt] {Start};
      \fill (V) circle (1.6pt) node[above right, inner sep=2pt] {$V$};
      \node[font=\scriptsize, anchor=west] at (5.55,0.3) {$Z$};
      \draw[densely dotted] (4.8,0.3) -- (5.55,0.3);
      \draw[pap arr, dashed] (S) .. controls (1.8,4.3) and (3.6,3.9) .. (V);
      \node[stepnr] at (2.5,3.95) {1};
      \node[font=\scriptsize, anchor=south] at (2.5,4.1) {\acs{ptp}};
      \draw[pap arr] (4.6,2.6) -- (4.6,1.05);
      \draw[pap arr] (5.0,1.05) -- (5.0,2.6);
      \node[stepnr] at (4.6,1.8) {2};
      \node[stepnr] at (5.0,1.8) {4};
      \node[stepnr] at (5.85,0.8) {3};
      \draw[densely dotted] (3.6,2.7) -- (4.7,2.7);
      \draw[densely dotted] (3.6,0.3) -- (4.13,0.3);
      \draw[{Stealth[length=1.5mm]}-{Stealth[length=1.5mm]}] (3.7,0.3) -- node[left] {$h$} (3.7,2.7);
    \end{tikzpicture}
  \end{minipage}\hfill
  \begin{minipage}[c]{0.54\textwidth}
    \small
    \newcommand{\stepmark}[1]{\tikz[baseline=(n.base)]\node[circle, draw, inner sep=0.8pt,
      font=\scriptsize\bfseries](n){#1};}
    \begin{tabularx}{\linewidth}{@{}l >{\raggedright\arraybackslash}X >{\raggedright\arraybackslash}X@{}}
      \toprule
      & \texttt{pick} & \texttt{place} \\
      \midrule
      vorab & $Z$ aus erkannter Pose, $V$ und $Z$ prüfen, Greifer öffnen
            & $Z$ = benannte Pose, $V$ und $Z$ prüfen \\
      \stepmark{1} & \multicolumn{2}{l}{\acs{ptp} nach $V$} \\
      \stepmark{2} & \multicolumn{2}{l}{\acs{lin} nach $Z$, langsam} \\
      \stepmark{3} & Greifer schließen, warten & Greifer öffnen, warten \\
      \stepmark{4} & \multicolumn{2}{l}{\acs{lin} zurück nach $V$} \\
      \bottomrule
    \end{tabularx}
  \end{minipage}
  \caption{Greif- und Ablegesequenz: Bewegung in der Seitenansicht (links) und Schritte von
  \texttt{pick} und \texttt{place} (rechts). Die Vorposition $V$ liegt um den
  Sicherheitsabstand $h$ über der Zielpose $Z$. Geprüft wird gegen die Arbeitsraumgrenzen.}
  \label{fig:m2-pick}
\end{figure}

Das Greifen an der Prüfposition nutzt dieselbe Sequenz. Statt einer erkannten Pose wird
dabei die fest hinterlegte Pose der Prüfposition übergeben, da das Bauteil dort in bekannter
Lage liegt.

\todo[inline]{Greifer-Rückmeldung auswerten: Nach dem Schließen prüfen, ob tatsächlich ein
Bauteil gegriffen wurde.}

% ------------------------------------------------------------------------------
\section{Feste Posen}
\label{sec:m2-posen}

Neben der erkannten Bauteilpose fährt der Roboter vier feste Posen an
(\cref{tab:m2-posen}). Sie werden am realen Aufbau im Handbetrieb angefahren (geteacht) und
in die Konfiguration übernommen. Die Ruheposition ist als benannter Punkt auf der Steuerung
gespeichert und wird direkt in Gelenkkoordinaten angefahren.

\begin{table}[htbp]
  \centering
  \caption{Feste Posen des Roboters}
  \label{tab:m2-posen}
  \begin{tabularx}{\textwidth}{@{}l X@{}}
    \toprule
    Pose & Zweck \\
    \midrule
    \texttt{home} & Ruheposition außerhalb der Kamerabilder; Start- und Endpunkt \\
    \texttt{inspection} & Ablage an der Prüfposition im Sichtfeld der Seitenkamera \\
    \texttt{bin\_good} & Ablage für Gutteile \\
    \texttt{bin\_bad} & Ablage für Schlechtteile \\
    \bottomrule
  \end{tabularx}
\end{table}

\todo[inline]{Posen am realen Roboter teachen und eintragen. Wiederholgenauigkeit an der
Prüfposition beurteilen, da die Sichtprüfung eine reproduzierbare Lage voraussetzt.}

% ------------------------------------------------------------------------------
\section{Kalibrierung zwischen Kamera und Roboter}
\label{sec:m2-kalibrierung}
\todo[inline]{Verfahren noch offen (siehe \cref{sec:koordinatensysteme}).}

% ------------------------------------------------------------------------------
\section{Sicherheitskonzept}
\label{sec:m2-sicherheit}

Die Anlage ist nicht sicherheitszertifiziert und wird nur unter Aufsicht betrieben. Die
Robotersteuerung ergänzt die Schutzfunktionen des \ac{cobot}s um mehrere
Softwaremaßnahmen:

\begin{description}
  \item[Reduzierte Geschwindigkeit] Im Testbetrieb sind die globale Geschwindigkeit der
    Steuerung und die Gelenkgeschwindigkeit auf \SI{20}{\percent} begrenzt. Die geradlinige
    Annäherung an das Bauteil erfolgt zusätzlich mit geringer Geschwindigkeit.
  \item[Arbeitsraumgrenzen] Vor jeder Bewegung wird geprüft, ob die Zielpose und die
    Vorposition innerhalb eines festgelegten Quaders im Roboter-\ac{ks} liegen. Liegt eine
    Pose außerhalb, wird die Bewegung nicht ausgeführt, und die Robotersteuerung meldet einen
    Fehler. Dadurch führt auch eine fehlerhafte Pose der Objekterkennung nicht zu einer
    Bewegung außerhalb des vorgesehenen Bereichs.
  \item[Stopp mit Quittierung] Ein Stopp beendet das Bewegungsprogramm auf der Steuerung.
    Danach werden alle Bewegungsbefehle abgelehnt, bis der Bediener den Stopp quittiert und
    der Roboter erneut vorbereitet wird (\cref{sec:m2-api}). Das Wiederanlaufen geschieht
    bewusst nicht automatisch.
  \item[Not-Halt] Der Hardware-Not-Halt des Roboters ist jederzeit in Reichweite.
\end{description}

% ------------------------------------------------------------------------------
\section{Inbetriebnahme}
\label{sec:m2-inbetriebnahme}

Für die schrittweise Inbetriebnahme am realen Roboter gibt es ein Prüfprogramm, das die
Funktionen einzeln ausführt: Verbindung und Statusabfrage, Anfahren der Ruheposition,
Öffnen und Schließen des Greifers, Anfahren einer festen Pose und schließlich ein
vollständiger Greif- und Ablegevorgang. Dieselben Schritte stehen in der Bedienoberfläche
zur Verfügung (\cref{ch:integration}). Vor dem Einsatz am Roboter wurde der gesamte
Programmpfad gegen die Nachbildung der Steuerung getestet (\cref{sec:teststrategie}).

\todo[inline]{IP-Adresse, Werkzeug und Betriebsartwechsel am realen Roboter verifizieren.}

% ------------------------------------------------------------------------------
\section{Zwischenfazit}
\todo[inline]{Nach den Versuchen ergänzen: Was hat funktioniert, welche Probleme traten auf?}
